\section{Introduction}
\label{sec:introduction}

Information lies at the heart of artificial intelligence.
From Shannon's entropy~\cite{shannon1948mathematical} guiding learning
algorithms to the Information Bottleneck~\cite{tishby2000information} shaping
representations, virtually every paradigm in AI can be viewed as
compressing, transferring, organizing, or extracting meaning from information.
Yet, while much of classical information theory is concerned with quantities
of uncertainty at a \emph{local} or \emph{pairwise} level,
\textbf{Structural Entropy (SE)}~\cite{Li2016} offers a fundamentally
distinct and complementary perspective: it captures the \emph{global
structure} and \emph{hierarchical uncertainty} embedded within complex
information systems.

At its core, SE provides a rigorous method to quantify and decode the
uncertainty embedded in the architecture of a complex system, modelled as a
graph or matrix.
Its significance for the AI community is twofold:

\begin{itemize}
  \item \textbf{Metric.}
        SE serves as a fundamental metric that quantifies a system's intrinsic
        uncertainty.
        Low SE signals a clear, efficient hierarchical organization---implying
        functional coherence and robustness.
        High SE indicates a system that is difficult to decompose, approaching
        the behaviour of a random graph.
        This analytical power is exploited in privacy-preserving graph analysis
        (deliberately increasing SE to obscure community
        structure~\cite{liu2019rem}), knowledge
        quantification~\cite{wang2023quantifyingKnowledge}, and AI safety
        evaluation~\cite{zeng2025si2af}.

  \item \textbf{Optimization principle.}
        SE provides a principled method to \emph{discover} a system's optimal
        hierarchical structure.
        Minimizing SE decodes the most efficient information-theoretic
        abstraction of a system, revealing a nested, hierarchical partitioning
        of its components---the ``true'' community
        structure~\cite{Li2016} or most effective decision
        hierarchy~\cite{zeng2025sidmJMLR}.
\end{itemize}

In the context of \emph{Transactions on Graph Intelligence and Network
Applications} (TGINA), SE is timely because modern graph intelligence systems
increasingly need more than accurate predictions: they need interpretable,
auditable, and reusable structural abstractions.
This is especially visible in graph-enhanced large language model pipelines,
including GraphRAG-style retrieval, knowledge-graph reasoning, tool-use
memory, and autonomous-agent interaction graphs.
These systems generate relational structures whose trustworthiness depends on
whether meaningful hierarchies can be separated from incidental links.
SE provides a principled way to measure such organization and, through
minimization, to construct candidate hierarchies for retrieval, reasoning, and
decision making.

\subsection*{Motivation: Why a Decade Survey Now?}

The foundational SE paper~\cite{Li2016} was published a decade ago.
In the intervening years, SE-based methods have proliferated across graph
pooling~\cite{wu2022pooling,wxb2024hi-part}, data
augmentation~\cite{wu2023sega,wang2023user}, graph structure
learning~\cite{zou2023segsl,duan2024structuralGSL}, reinforcement
learning~\cite{zeng2025sidmJMLR,zeng2023hierarchicalSISA,zeng2024effectiveSI2E},
bioinformatics~\cite{zhang2021supertad,chen2022pyseat,li2023dedoc2},
social-network analysis~\cite{peng2024zjyun,yang2024sebot,cao2024socialevent},
and pattern recognition~\cite{zhu2023HiTIN,zeng2023lesion,wang2024secodecSpeech}.
Despite this rapid proliferation, a comprehensive and systematic review
consolidating these advances has been lacking.

This article aims to fill that critical gap.
The present work substantially extends our earlier conference
survey~\cite{su2025ijcai} published at IJCAI 2025:

\begin{enumerate}
  \item \textbf{Expanded theory.}
        We provide a self-contained, rigorous treatment of SE's mathematical
        foundations, including proofs of key properties and formal comparisons
        with alternative graph-complexity measures (modularity, map equation,
        von Neumann entropy).

  \item \textbf{Comprehensive algorithm catalogue.}
        We survey all known SE minimization algorithms, extending the
        conference table with complexity analysis, constraint categories, and
        reproducibility notes.

  \item \textbf{Reproducibility benchmark.}
        We present a new empirical evaluation comparing SE-based graph
        clustering objectives against Louvain, Infomap, and differentiable
        baselines across multiple datasets and random seeds using the
        \texttt{glass-jax} companion library.

  \item \textbf{Expanded application coverage.}
        We include additional application domains (speech processing,
        geoscience knowledge quantification, LLM integration) and updated
        literature through 2026.

  \item \textbf{Open problems.}
        We provide a structured discussion of open research challenges and
        future directions that were absent from the conference version.
\end{enumerate}

\subsection*{Paper Organization}

The remainder of this survey is organized as follows.
Section~\ref{sec:related} positions SE against modularity, flow-based,
spectral, and deep graph-clustering methods.
Section~\ref{sec:taxonomy} introduces a taxonomy of SE research and its
formal definitions.
Section~\ref{sec:theory} covers theoretical foundations.
Section~\ref{sec:algorithms} surveys computational algorithms for SE
minimization.
Section~\ref{sec:learning} examines SE-driven learning paradigms (graph
learning and reinforcement learning).
Section~\ref{sec:applications} reviews cross-domain applications.
Section~\ref{sec:benchmark} presents the reproducibility benchmark.
Section~\ref{sec:open} discusses open problems and future directions.
Section~\ref{sec:conclusion} concludes.
